\documentclass{article}
\usepackage{iclr2027_conference,times}
\input{math_commands.tex}
\usepackage{hyperref}
\usepackage{url}
\usepackage{booktabs}
\usepackage{array}
\hypersetup{hidelinks,pdfauthor={Anonymous authors},pdfsubject={Prospective study; confirmatory results pending}}

% Override from the command line; all 12 x 3 x 3 combinations are supported.
\providecommand{\TitleVersion}{1}
\providecommand{\AbstractVersion}{1}
\providecommand{\IntroVersion}{1}
\newcommand{\PBPFSelectorError}[2]{\PackageError{PBPF}{Invalid #1 selector}{Use an integer from 1 to #2.}}
\ifnum\TitleVersion<1 \PBPFSelectorError{title}{12}\fi
\ifnum\TitleVersion>12 \PBPFSelectorError{title}{12}\fi
\ifnum\AbstractVersion<1 \PBPFSelectorError{abstract}{3}\fi
\ifnum\AbstractVersion>3 \PBPFSelectorError{abstract}{3}\fi
\ifnum\IntroVersion<1 \PBPFSelectorError{introduction}{3}\fi
\ifnum\IntroVersion>3 \PBPFSelectorError{introduction}{3}\fi
% Normalize equivalent integer spellings before macro/file-name lookup.
\edef\TitleVersion{\number\numexpr\TitleVersion\relax}
\edef\AbstractVersion{\number\numexpr\AbstractVersion\relax}
\edef\IntroVersion{\number\numexpr\IntroVersion\relax}
% Exactly twelve candidates; title 1 is recommended.
\newcommand{\PBPFTitleChoice}[2]{\expandafter\def\csname PBPFTitle#1\endcsname{#2}}
\PBPFTitleChoice{1}{A Failed Test Is Not a Diagnosis: Particle-Belief Program Filtering for Code Repair}
\PBPFTitleChoice{2}{Failure Is Partially Observed: Particle Beliefs for Execution-Grounded Code Repair}
\PBPFTitleChoice{3}{PBPF: Sequential Particle Beliefs for Coherent Code Repair}
\PBPFTitleChoice{4}{Tests as Evidence: Particle Filtering Hidden Failure Modes for Program Repair}
\PBPFTitleChoice{5}{Belief Before Repair: Sequential Bayesian Diagnosis for Code Language Models}
\PBPFTitleChoice{6}{From Test Failures to Failure Beliefs: Particle-Conditioned Code Repair}
\PBPFTitleChoice{7}{Diagnose, Then Repair: Execution-Conditioned Particle Beliefs for Code LLMs}
\PBPFTitleChoice{8}{One Posterior, One Repair: Coherent Particle Conditioning for Iterative Code Generation}
\PBPFTitleChoice{9}{Beyond Pass Rates: Latent Failure Beliefs for Iterative Code Repair}
\PBPFTitleChoice{10}{Coherent Repairs from Ambiguous Tests: Particle-Belief Program Filtering}
\PBPFTitleChoice{11}{Which Bug Caused the Failure? Sequential Particle Beliefs for Test-Time Code Repair}
\PBPFTitleChoice{12}{Maintaining Many Diagnoses, Sampling One Repair: Particle Beliefs for Code LLMs}

% Deliberately no path for silently importing empirical or legacy results.
\newcommand{\PendingResult}[1]{\textsf{[PENDING: #1]}}
\newcommand{\ResultFiniteAudit}{\PendingResult{finite audit}}
\newcommand{\ResultPrediction}{\PendingResult{future-outcome NLL}}
\newcommand{\ResultCalibration}{\PendingResult{Brier and calibration}}
\newcommand{\ResultRepair}{\PendingResult{hidden Pass@1}}
\newcommand{\ResultCoherence}{\PendingResult{coherence controls}}
\newcommand{\ResultReplication}{\PendingResult{locked replication}}
\newcommand{\ResultResources}{\PendingResult{resource accounting}}

\title{\raggedright\csname PBPFTitle\TitleVersion\endcsname}
\author{Anonymous authors}
% \iclrfinalcopy % Never enable for anonymous submission.

\begin{document}
\maketitle
\begin{abstract}
\input{variants/abstract_\AbstractVersion}
\end{abstract}
\input{variants/intro_\IntroVersion}

\section{Candidate-specific particle beliefs}
\label{sec:method}
Let $x$ be a public task, $c_r$ a candidate version in one repair lineage,
$\Delta_r$ its patch from the selected parent, and
$E_r=(e_{r,1},\ldots,e_{r,J})$ its ordered visible outcomes. A diagnostic
particle $z_r^m$ is not an independently decoded candidate program.
The model factorizes along a lineage as
\begin{equation}
 p(z_0\mid x,c_0)p(E_0\mid z_0,x,c_0)
 \prod_{r\geq1}p_\psi(z_r\mid z_{r-1},x,c_r,\Delta_r)
 p_\theta(E_r\mid z_r,x,c_r).
\end{equation}
The latent is fixed within each version. We factor the observation likelihood
over ordered tests into categorical probabilities for \texttt{PASS},
\texttt{WRONG\_OUTPUT}, \texttt{RUNTIME\_EXCEPTION}, \texttt{TIMEOUT}, and
\texttt{COMPILE\_ERROR}. Infrastructure failure is recorded separately, not
learned as a sixth outcome. This conditional-independence model is an
approximation; correlated tests and misspecified likelihoods remain limitations.

\paragraph{Creation and observation are different operations.}
At a root, draw $z_0^m\sim q_{\rm root}(z\mid x,c_0,e_{0,1})$ and assign
unnormalized log mass
$-\log P+\log p(z_0^m\mid x,c_0)+\log p_\theta(e_{0,1}\mid z_0^m,x,c_0)
-\log q_{\rm root}(z_0^m\mid x,c_0,e_{0,1})$.
At a patch boundary, select an ancestor $a_m$ explicitly and sample
$z_r^m\sim q_\phi(z\mid z_{r-1}^{a_m},x,c_r,\Delta_r,e_{r,1})$.
The first child log mass is
\begin{align}
 \ell_r^m={}&b_m+\log p_\psi(z_r^m\mid z_{r-1}^{a_m},x,c_r,\Delta_r)
 +\log p_\theta(e_{r,1}\mid z_r^m,x,c_r)\nonumber\\
 &-\log q_\phi(z_r^m\mid z_{r-1}^{a_m},x,c_r,\Delta_r,e_{r,1}).
 \label{eq:child}
\end{align}
For deterministic enumeration of all parent indices, $b_m=\log w_{r-1}^{a_m}$.
For sampled ancestors $a_m\sim\rho$, including repeated draws,
$b_m=-\log P+\log w_{r-1}^{a_m}-\log\rho(a_m)$, with proposal support covering
the parent target. In particular, $\rho=w_{r-1}$ gives uniform base masses;
parent weights must not be counted twice. The incremental log normalizer is
$\log\sum_m\exp\ell_r^m$, before normalization or resampling. We never
renormalize the sampled-ancestor base masses before this calculation.

For subsequent tests on unchanged code, the update is likelihood-only:
$\ell_{r,j}^m=\log w_{r,j-1}^m+\log p_\theta(e_{r,j}\mid z_r^m,x,c_r,t_j)$.
Transition and proposal terms require an actual newly sampled latent and are
evaluated at that sample. Candidate hashes key independent states; siblings
inherit only their selected parent. When $\mathrm{ESS}<P/2$, systematic
resampling is followed by one Metropolis-adjusted Langevin move that includes
forward and reverse proposal densities~\citep{gilks2001}. Uncorrected jitter
is a named biased ablation, not the primary filter.

\paragraph{Learning and calibration.}
The primary model uses a 32-dimensional latent and eight particles. Frozen
actor hidden states are attention-mask-mean-pooled for semantic features.
Learned diagonal-Gaussian root, transition, and child-proposal heads accompany
a five-way likelihood. The belief objective combines an SMC evidence objective
with a proper future-prediction score~\citep{maddison2017,naesseth2018}:
\begin{equation}
 \mathcal L_B=-\mathbb E[\log\widehat Z_{\rm SMC}]
 +\lambda_{\rm future}\mathcal L_{\rm future},\qquad
 \mathcal L_{\rm future}=-\sum_{j>k}\log\sum_m w_k^m
 p_\theta(e_j\mid z^m,x,c,t_j).
\end{equation}
Training samples prefixes $k\in\{1,2,4\}$; the primary endpoint fixes $k=4$.
Discrete ancestor choices are detached. The categorical NLL is unweighted;
one scalar temperature fitted on source-disjoint calibration data is reused
in filtering and reported predictions~\citep{guo2017}.

\section{Coherent posterior-conditioned repair}
\label{sec:action}
An eight-token learned soft prefix maps one latent to the frozen actor's input
embeddings~\citep{li2021prefix}. The belief is frozen before training this
projector. For a target repair $a^*$, the objective is
\begin{equation}
 \mathcal L_A=-\log\sum_m\operatorname{stopgrad}(w_m)
 \exp\left\{\sum_t\log\pi(a_t^*\mid a_{<t}^*,x,c,E,z^m)\right\}.
 \label{eq:mixture}
\end{equation}
Component responsibilities are first computed without gradients; each
component graph is then replayed serially for bounded-memory gradient
accumulation. At inference, one component is held fixed for all tokens of a
repair. Across multiple calls, systematic component sampling allocates
posterior mass while preserving within-call coherence. A full decode for each
particle would consume additional actor work and is not our primary method.
Token-wise remix is explicitly a fault ablation: in general,
$\sum_m w_m\prod_t\pi_m(a_t\mid a_{<t})$ differs from
$\prod_t\sum_m w_m\pi_m(a_t\mid a_{<t})$.

\section{Experimental Setup}
\label{sec:setup}
The following protocol is prospective. Configuration values specify intended
experiments, not measurements or completed runs. The Slurm dispatcher and
production-factory interface are implemented, but a complete production
train/select/calibrate/evaluate factory is not shipped. An operator must
integrate and audit the scientific handlers and live evaluator before launch.

\paragraph{Research questions and stages.}
RQ1 asks whether particle inference matches exact finite posteriors. RQ2 asks
whether a candidate-specific belief predicts withheld outcomes beyond matched
deterministic memories. RQ3 asks whether that predictive value improves
equal-budget hidden repair. RQ4 tests continuation coherence, and RQ5 tests
robustness to evidence, model, and benchmark changes. S0 checks integrity on
16 tasks per dataset; S1 audits exact finite inference; S2 trains and freezes
beliefs and tests prediction; S3 tests repair only after the predictive gate;
S4 evaluates locked replications. Synthetic offline execution is smoke-only
and cannot establish any empirical claim about code actors.

\paragraph{Data and information partitions.}
The finite audit specifies 50,000 training, 5,000 development, and 5,000 test
cases with disjoint latent families and generator templates, ten tests per
case, four visible and six future. Its generator has 120 families, ten
templates per family, 500 cases per family, and 32 discrete latent states;
Dirichlet priors and five-way likelihood tables permit exact posterior
calculation. These are generator-contract counts, not a claim that cases have
already been generated.

PBPF-RBR adapts RunBugRun~\citep{runbugrun}: eligible Python records have at
least ten active tests, a fixed version passing all tests, a buggy version
failing at least one, and code length at most 2,048 actor tokens. Official
boundaries are preserved and connected leakage components based on problem
identity and normalized code/test signatures are removed across splits.
Deterministic caps are 12,000/2,000/2,000 for official training/development/test.
Training reserves at most 1,000 source-disjoint calibration tasks, leaving at
most 11,000 for fitting; development is split into tune/select halves of at
most 1,000 each. SHA-256 ordering selects components without breaking them,
so actual counts may be below caps. Tests are SHA-256 ordered, with the first
four visible. Gold fixes and evaluator-only expected outputs never enter
generator prompts or test-time selection.

CodeARC-Replay uses the anonymous CodeARC variant~\citep{codearc}, fixed
invocation indices 0--3 as visible and 4--9 as evaluator-only. It is a
\emph{replay} protocol, not the official interactive leaderboard; the target
program remains evaluator-only. PBPF-EvalPlus uses all 164 HumanEval+ and
378 MBPP+ tasks~\citep{evalplus}. Base tests supply visible feedback and
exactly deduplicated plus-only tests are hidden, run only after final hashes
are sealed. The four-test prefix is the primary prediction endpoint and the
RBR/CodeARC visible limit, not permission to relabel EvalPlus base cases as
hidden. This is a repair protocol, not a zero-shot leaderboard claim.
The locked LiveCodeBench replication uses the release-v6 tasks absent from
release-v5~\citep{livecodebench}; its size and CodeARC's size are determined by
the pinned snapshot inventory, not guessed. SWE-bench is outside the core claim.

\paragraph{Models and inference.}
Qwen2.5-Coder-7B-Instruct is primary. Qwen3-8B with thinking disabled and
DeepSeek-Coder-6.7B-Instruct are the additional model replications; the two Qwen
models are not described as independent model families. The 1.5B Qwen model is
opt-in plumbing only. Each model and tokenizer uses its separately pinned full
immutable repository revision and native chat template. Inference uses BF16,
context length 8,192, output cap 1,024 tokens, temperature
0.8, and top-$p$ 0.95. Primary conditioning uses $P=8$, $d_z=32$, eight soft
prefix tokens, systematic resampling at $\mathrm{ESS}<P/2$, and one corrected
MALA move. Low-rank K, V, K+V injection and QLoRA are secondary ablations or
replications, not the headline model.

\paragraph{Training and development locks.}
Frozen actor last-layer features are attention-mask-mean-pooled and detached.
Stage B uses AdamW, weight decay 0.01, gradient clipping 1.0, learning rates
$\{10^{-4},3\mathord\times10^{-4}\}$, effective batches $\{16,32\}$,
maximum steps $\{2000,5000\}$, $\lambda_{\rm future}\in\{0.5,1.0\}$, and
MALA step sizes $\{0.01,0.05,0.1\}$. The tune split ranks task-macro prefix-four
NLL, Brier, fewer steps, smaller batch, smaller learning rate, smaller future
weight, then smaller MALA step. The actor projector uses AdamW with the same
learning-rate, batch, and maximum-step grids. Beliefs are frozen before actor
conditioning; trainable module counts and actual optimization work must be
reported. One scalar temperature is fitted once on calibration data.

The deterministic comparator is selected on development-select from full
transcript, window-two, window-four, DeepSets, matched GRU, and scalar
correctness, ranking prefix-four NLL, Brier, then lexical arm name. The two
replication baselines rank development visible-selected Pass@1, fewer actor
output tokens, then lexical name. Choices are locked before test prediction,
test repair, and replication. Selected hyperparameters are refitted on training
only; neither development-select, calibration, nor test enters parameter
fitting. Model-specific modules are retrained for each replication under the
locked hyperparameters and calibrated once without locked-test tuning.

\paragraph{Baselines and budgets.}
Prediction compares prior, last outcome, full transcript, windows two/four,
DeepSets, matched GRU, scalar correctness, MAP, posterior mean, PBPF, and five
controls: shuffled evidence, wrong candidate, masked outcomes, random latent,
and faulty shared belief. Deterministic parameter matching includes all unique
learned Stage-B parameters, including frozen learned components, encoders,
projections, biases, and the five-way head, within 5\%; it excludes the shared
frozen actor/semantic encoder, particle values, buffers, optimizer, and the
Stage-C projector. Width ties minimize relative mismatch then width.

Repair compares independent one-shot repairs, a controlled full-transcript
Self-Debugging adaptation~\citep{selfdebug}, paper-spec REx~\citep{rex}, the
locked strongest deterministic belief, a paper-spec Rollout Roulette
adaptation~\citep{roulette}, and PBPF sample-once. These local adaptations are
not claims of bitwise official reproduction. Self-Debugging explanations must
fit inside the four allowed decodes. REx uses visible pass-fraction candidate
allocation. Rollout Roulette resamples four ordered partial trajectories
before completion under a two-segment profile; partial and discarded prefix
tokens are metered separately. A no-repair selector has lower compute and is
not included as an equal-repair-budget baseline. LDB~\citep{ldb} and other
trace-rich methods belong only to an enhanced-information table. RLEF is
related actor-training work~\citep{rlef}, not an unimplemented matched arm.

For each task and seed, eight genuine actor samples are generated once and
their exact hashes shared across arms. Every matched repair arm receives
exactly four repair decodes, the same ordered visible tests, status-only
feedback, output limits, and no early stopping. Independent repair uses four
distinct initial slots without later repair feedback. The initial selected
Pass@1, hidden pass@8 oracle ceiling, and final hidden Pass@1 are distinct
metrics. Exactly one final program is chosen by visible pass fraction with
fixed tie-breaking and sealed before hidden evaluation. Belief calls are not
actor decodes, but every full particle-conditioned decode increases effective
decode multiplicity and must be charged as such.

\paragraph{Ablations and controls.}
Development uses 256 RBR tasks and three seeds, one factor at a time, not a
Cartesian ablation sweep. Factors cover $P\in\{1,2,4,8,16,32\}$; learned versus
uniform proposal, likelihood, and transition; proposal correction on/off;
resampling never/every step/ESS; corrected move/off/biased jitter; evidence
last/window/full/orderless; sample-once/MAP/mean/token-remix/full ensemble;
no-op/random/text/soft-prefix/K/V/KV injection; and prefix lengths 2/4/8/16.
The confirmatory configuration and five causal controls are locked. Shuffling
permutes observed outcome/test pairings, wrong-candidate feedback uses a
same-task same-round donor or masked fallback, masking removes outcome-derived
features, random latents match posterior weighted RMS radius, and shared belief
is labeled faulty. Full traces are a different information condition.

\paragraph{Metrics, statistical units, and gates.}
Formal seeds are 1701, 1702, and 1703. The independent unit is a source-problem
cluster. Metrics aggregate tests within candidate, candidates within task,
then seeds within task. Inference uses 10,000 paired stratified cluster-bootstrap
replicates, seed 314159, shared resample indices across arms, strata for dataset,
official split, and source language, and 95\% intervals. Secondary comparisons
use Holm correction; cross-model summaries have fixed equal weights and every
model-by-dataset cell remains visible.

S1 compares exact Bayes with $P\in\{1,4,8,16,32\}$: at 32 particles, median
posterior KL must be at most 0.05 nat, future NLL within 2\% of exact, and
randomized 90\% HPD coverage in $[0.87,0.93]$. Gate B requires at least 5\%
task-macro prefix-four future-NLL improvement over the development-selected
deterministic comparator and a positive paired 95\% lower bound on log-relative
improvement. A one-sided Brier non-inferiority check uses absolute margin 0.01
on the multiclass $[0,2]$ scale, a prospective tolerance rather than an empirical
estimate. Shuffled, wrong-candidate, and random controls must each remove at
least 80\% of the gain. Failure stops confirmatory repair. Gate C requires at
least three absolute percentage points in final hidden Pass@1 over both
Self-Debugging and the strongest deterministic arm, with both paired lower
bounds positive and Gate B retained; failure blocks locked confirmatory
replication. S4 requires agreeing directions for both additional models and
a positive fixed-equal-weight cross-model lower bound. Main-table eligibility
requires Gates B and C and the factory-produced locked replication gate.

S2 evaluates the sealed RBR test and CodeARC-Replay. S3 runs primary-model repair
on RBR, full EvalPlus, and CodeARC-Replay. S4 runs PBPF plus the two preselected
baselines on full EvalPlus for Qwen3 and DeepSeek, and on CodeARC-Replay plus
the locked LiveCodeBench slice for Qwen2.5. No model-by-dataset cell is replaced
by a more favorable pooled result.

\paragraph{Isolation, provenance, and resource accounting.}
Generator and evaluator are separate processes with separate manifests. The
generator cannot access hidden test source, expected hidden outputs, gold code,
gold patches, or final hidden outcomes. Dataset/model/tokenizer/prompt revisions,
materialized snapshot and weight hashes, native chat-template hashes, split
seeds, container digests, test order, and clusters are bound to the run identity;
resume rejects fingerprint mismatches. Stable SHA-256 sharding and private
atomic shard outputs permit deterministic merge with duplicate, missing-key,
and checksum rejection. The intended entry point is
\texttt{bash scripts/run\_iclr.sh} with config
\texttt{configs/experiments/iclr\_pbpf.yaml}, profile
\texttt{slurm\_h200x16}, and \texttt{--resume}. Sixteen independent single-GPU
workers are a launch profile, not a claimed resource measurement.

Formal evaluator authorization additionally requires an independently launched,
distinct-UID worker, a pre-provisioned evaluator-owned trust anchor, an exact
signed seal digest, pinned roots, and an external sandbox specification.
Same-UID and stdio-only evaluation is permanently nonconfirmatory. Stage-B
future scoring has its own immutable prediction seal; final repair evaluation
uses a separate exact selection mapping. Accepted receipts are unique per
sealed work ID, but a crash can cause another physical attempt, so failed and
retried actor/verifier work is reserved and charged rather than hidden.

Before scaling, a disjoint 64-task resource pilot must measure throughput and
abort for a projected-cost deviation above 25\%, memory above 90\%, or
infrastructure failures above 1\%. Ledgers separately report generation calls,
belief forwards, verifier cases/suites, input/output tokens, CPU seconds, GPU
hours, wall time, peak allocated memory, and disk bytes. Static actor-token and
verifier-suite bounds cover one three-seed evaluation sweep per cell, not
training, grid search, or retries. Full-decode equivalents and invocation
counts are distinct, including partial-prefix requests. Full-workflow costs
require operator-approved GPU/CPU budgets and the resource pilot.
Forced execution after
a failed gate is explicitly nonconfirmatory and excluded from main tables.


\section{Results status and interpretation}
\label{sec:results}
\textbf{This is a prospective manuscript, not a completed empirical report.}
No formal learned-model prediction, hidden-test repair, or cross-model
replication has been established by this package. Software tests and synthetic
smoke outputs do not establish those claims. The reported result slots are
\ResultFiniteAudit, \ResultPrediction, \ResultCalibration, \ResultRepair,
\ResultCoherence, \ResultReplication, and \ResultResources.
If prediction improves without repair gains, the belief may be informative
but poorly coupled to the actor. If repair improves without the predictive
gate or without sensitivity to evidence interventions, that outcome will not
support the diagnosis claim. Failed gates and forced downstream runs must be
reported as such; no pending entry may be replaced from legacy diagnostics.

\section{Limitations}
A learned latent need not correspond to a unique human-interpretable bug, and
predictive calibration does not identify causal mechanisms. Finite-particle
approximation, conditional-independence assumptions, and proposal misspecification
can invalidate a useful posterior interpretation. A frozen actor may ignore
or misinterpret a prefix, and the method still incurs substantial training
and inference costs. Public benchmarks may have been seen during pretraining;
the time-sliced replication reduces but cannot prove the absence of
contamination. Replay and repair protocols are not interchangeable with
official interactive or zero-shot leaderboards. Results, once available,
will not establish applicability to arbitrary repositories, languages,
security-critical software, or unconstrained execution environments.
\label{sec:main-end}

\section*{AI use statement}
Generative AI assistance, including ChatGPT/Codex, was used to propose and
refine the method and experimental design, develop mathematical formulations
and derivations, implement and test software, research and summarize related
literature, format references, and draft and revise this manuscript and its
artifact documentation. This includes assistance with data-preparation code
and the design of synthetic audit generators; it does not mean that the
formal datasets or confirmatory experiments have been executed. No
confirmatory measurements were generated or inferred for the result slots.
This statement does not assert that human review of all AI-assisted material
has been completed. Verification records describe the checks actually run;
the authors remain responsible for validating the content and completing
the required submission disclosure before submission.

\section*{Reproducibility statement}
Sections~\ref{sec:method}--\ref{sec:setup} specify the probability model,
objectives, information partitions, selection rules, and prospective gates.
The source package includes switchable narrative variants, unmodified
conference assets, build instructions, and provenance notes. The associated
implementation is intended to emit immutable configuration, data, candidate,
and budget manifests. A buildable paper or a successful offline smoke test
is not evidence that the external formal backend, dataset preparation,
training, or GPU experiments have completed. Those artifacts and their
checks must accompany any future numerical revision.

\section*{Ethics statement}
Generated programs require isolated execution with restricted privileges,
resource limits, and no unintended network or filesystem access. Passing
benchmark tests does not certify code safety or correctness. Dataset access,
redistribution, and licensing must be checked per upstream source; we do not
infer permission to redistribute benchmark payloads from a repository's
software license. The evaluator firewall and explicit pending results are
intended to reduce leakage and unsupported scientific claims.

\bibliography{references}
\bibliographystyle{iclr2027_conference}
\appendix
\section{Additional protocol details}
\label{app:details}
The packaged protocol lock binds the reviewed configuration files and the
approved design by SHA-256. It is not a claim of a public preregistration and is not
evidence that the configured immutable revisions resolved successfully.
Execution manifests must independently verify the revisions actually loaded.

Final selection uses visible pass fraction only, then lower repair round,
lower initial slot, lexicographic source hash, and lexicographic version ID.
Source-component caps are upper bounds: components are never split merely to
hit a nominal count, and actual inclusions and exclusion reasons must be
reported. Main tables must show every model-by-dataset cell, not only a
pooled estimate. Any empirical revision must attach the source artifact hash,
selection lock, denominator, seed set, confidence interval, and budget ledger
for each numerical claim.
\end{document}
